\section{Related Work}
\label{sec:related-work}

\paragraph{Reference-signal operators.}
Regression-based ocular correction estimates how eye potentials propagate to each scalp channel and subtracts the predicted contribution \cite{gratton1983ocular,semlitsch1986ocular,schlogl2007regression}; its per-subject ``propagation factors'' are the ancestor of our calibrated matrix, and later work showed that eye-related subspaces estimated once from a calibration recording can be applied on-line with matrix operations \cite{kobler2020sgeyesub}. The lineage is active: ARMBR estimates a participant-specific blink filter from seconds of data \cite{alkhoury2025armbr}, and adaptive filters track the propagation during use \cite{he2004rls}. This literature established both the value of a calibration segment and the linear propagation model we build on, but it treats clean EEG only implicitly, as whatever remains after subtraction. Our method keeps the calibrated propagation model of this lineage and adds what it lacks: a learned distribution of clean EEG that decides how the subtraction should be corrected, a shrinkage estimator that protects a short calibration (Section~\ref{sec:res-paired}), and a reliability criterion with a measured fail-safe.

\paragraph{Reference-free subspace methods.}
Blind source separation with automated component selection (ICA with ICLabel \cite{jung2000bss,pion2019iclabel}) and artifact subspace reconstruction (ASR) \cite{mullen2015asr,chang2020asr} need no EOG channel. ASR in particular normalizes our design premise: it also requires per-subject calibration data, and its documented sensitivity to calibration quality is exactly the failure mode our explicit reliability rule addresses --- ASR fails silently on a poor calibration, whereas our system withholds its guide and says so. We evaluate both families as anchor rows on our protocol and find that they buy attenuation with retention (Section~\ref{sec:res-natural}).

\paragraph{Population denoising networks and diffusion.}
Supervised networks trained on semi-simulated corpora such as EEGdenoiseNet learn flexible mappings from contaminated to clean signals \cite{zhang2021eegdenoisenet,sun2020rescnn,yu2022deepseparator,chuang2022icunet,lai2024cleegn}, calibration-free models extend this to unseen montages \cite{marcosmartinez2025eegoarnet}, and reconstruction accuracy on the shared benchmark has largely saturated under architecture iteration \cite{bindra2026capacity}, with the field's evaluation practice shifting toward real-data and task-aware endpoints \cite{xiang2025task}. Diffusion models learn to invert a prescribed noising process \cite{ho2020ddpm,song2021scoresde} and, conditioned on a degraded observation, restore images, audio, and biosignals \cite{saharia2022palette,richter2023storm,li2023descod}; EEGDfus and D4PM apply conditional diffusion to EEG artifact removal \cite{huang2025eegdfus,shao2025d4pm}. In all of these the participant variation is absorbed into shared parameters: none conditions on a recorded reference channel, adapts per subject, or reports uncertainty. The closest prior work to ours in intent is the domain-specific denoising diffusion model of \citet{duan2023dsddpm}, which makes diffusion-based EEG denoising subject-aware by learning subject-specific noise distributions inside the network, so that the subjects to be denoised participate in training. We share the goal and the diffusion backbone but place the subject information elsewhere: the network stays subject-independent, and subject specificity enters through a closed-form propagation estimate from the evaluated user's own calibration segment. The two designs answer different deployment questions --- learning subject structure into the weights suits a fixed user pool; measuring it at calibration time extends to users the network has never seen, which is the setting we evaluate. Our supporting experiment (Section~\ref{sec:res-representation}) suggests the two views reconcile: calibrated subject conditioning also works through a learned interface, and what is essential is that the conditioning is fitted to the target user rather than supplied as a training-time identity.

\paragraph{Guided restoration with measured operators.}
When the degradation operator is known, restoration methods such as DDRM and diffusion posterior sampling use it explicitly during sampling \cite{kawar2022ddrm,chung2023dps}; when it is unknown, blind approaches estimate operator and signal jointly \cite{chung2023blinddps,murata2023gibbsddrm}. Our setting sits between these cases: the operator is unknown a priori but \emph{measurable} --- a calibration segment provides it in closed form before sampling begins, with an uncertainty that we propagate into the sample distribution. The statistical device we use to stabilize the measurement, empirical-Bayes shrinkage of per-subject parameters toward a population prior, has precedent in cross-subject transfer for BCI decoders \cite{jayaram2016transfer}; to our knowledge it had not been applied to artifact operators.

\paragraph{Uncertainty for denoised signals.}
Probabilistic diffusion reconstruction of time series conventionally reports CRPS from posterior samples \cite{tashiro2021csdi,alcaraz2023sssd}, and the image-restoration literature has developed calibrated per-pixel intervals with distribution-free guarantees \cite{angelopoulos2022imageimage}. In EEG artifact removal, the only uncertainty-aware method we are aware of uses uncertainty internally as a training signal, without intervals or coverage \cite{jin2023udnet}; diffusion denoisers that draw multiple samples average them and discard the spread \cite{li2023descod}. Because diffusion posterior samples are not automatically calibrated \cite{crafts2025bipsda}, we evaluate our intervals by empirical coverage and the continuous ranked probability score \cite{gneiting2007strictly}, at frozen settings, on held-out users --- to our knowledge the first coverage-evaluated predictive intervals in EEG artifact removal.
